% !TEX root = ../algebra_02.tex

\section{Внутренняя прямая сумма линейных подпространств, эквивалентные определения}

\begin{Definition}{Внутренняя прямая сумма}{}
	Говорят, что линейное пространство $V$ \emph{раскладывается во внутреннюю прямую сумму} своих подпространств $W_1, \ldots, W_k \leq V$, и пишут $V = W_1 \oplus \cdots \oplus W_k$, если для каждого вектора $v \in V$ существует единственное представление в виде суммы:
	\[
	v = w_1 + \cdots + w_k, \quad \text{где } w_i \in W_i.
	\]
\end{Definition}

Существуют удобные критерии, позволяющие проверить, является ли сумма подпространств прямой, не проверяя единственность разложения для всех возможных векторов.

\begin{Theorem}{Критерий для двух слагаемых}{}
	Для подпространств $W_1, W_2 \leq V$:
	\[
	V = W_1 \oplus W_2 \iff V = W_1 + W_2 \text{ и } W_1 \cap W_2 = \{0\}.
	\]
\end{Theorem}

\begin{proof}
	$(\Rightarrow)$: Предположим, $V = W_1 \oplus W_2$. Равенство сумме очевидно по определению. Если $w \in W_1 \cap W_2$, то нулевой вектор можно представить двумя способами:
	$0 = 0 + 0$ (где $0 \in W_1$, $0 \in W_2$) и $0 = w + (-w)$ (где $w \in W_1$, $-w \in W_2$). Из-за единственности представления должно выполняться $w = 0$, откуда $W_1 \cap W_2 = \{0\}$.
	
	$(\Leftarrow)$: Пусть $V = W_1 + W_2$ и пересечение тривиально. Допустим, вектор $v$ имеет два представления: $v = w_1 + w_2 = w_1' + w_2'$. 
	Тогда перенесем слагаемые: $w_1 - w_1' = w_2' - w_2$.
	Вектор слева принадлежит $W_1$, а вектор справа --- $W_2$. Следовательно, оба они лежат в пересечении $W_1 \cap W_2 = \{0\}$. Это значит, что $w_1 - w_1' = 0$ и $w_2' - w_2 = 0$, то есть $w_1 = w_1'$ и $w_2 = w_2'$. Представление единственно.
\end{proof}

Для произвольного числа слагаемых критерий обобщается следующим образом (недостаточно попарного тривиального пересечения!).

\begin{Theorem}{Критерий для нескольких слагаемых}{}
	Пространство $V$ раскладывается в прямую сумму $V = W_1 \oplus \cdots \oplus W_k$ тогда и только тогда, когда выполнены два условия:
	\begin{enumerate}
		\item $W_1 + \cdots + W_k = V$;
		\item для каждого $j = 1, \dots, k$ выполнено: 
		\[
		W_j \cap (W_1 + \cdots + \widehat{W_j} + \cdots + W_k) = \{0\},
		\]
		(крышка означает пропуск слагаемого $W_j$).
	\end{enumerate}
\end{Theorem}
